% Generated from locale/tr/content/computability/recursive-functions/general-recursive-functions.tex; do not hand-edit.


\olfileid[tr]{cmp}{rec}{gen}
\olsection{Genel Özyinelemeli Fonksiyonlar}

Tümel fonksiyonlardan oluşan bir küme elde etmenin başka bir yolu vardır.
Her doğal sayı dizisi $\vec z$ için $f(x,\vec z)=0$ olacak bir $x$ varsa
tümel $f(x,\vec z)$ fonksiyonuna \emph{düzenli} diyelim. Başka bir deyişle
düzenli fonksiyonlar, sınırsız arama uygulandığında tümel bir fonksiyon elde
edilen fonksiyonların tam kendileridir. Sınırsız aramayı tutucu biçimde
düzenli fonksiyonlarla sınırlayabiliriz.

\begin{defn}
\ollabel{defn:general-recursive}
\emph{Genel özyinelemeli fonksiyonlar} kümesi, doğal sayıların çeşitli
kuvvetlerinden doğal sayılara giden fonksiyonlar arasında sıfır, ardıl ve
izdüşümleri içeren ve bileşke, ilkel özyineleme ile \emph{düzenli}
fonksiyonlara uygulanan sınırsız arama altında kapalı olan en küçük kümedir.
\end{defn}

Her genel özyinelemeli fonksiyon açıkça tümeldir.
\olref{defn:general-recursive} ile \olref[par]{defn:recursive-fn} arasındaki
fark, ikincisinde süreç boyunca kısmi özyinelemeli fonksiyonların
kullanılabilmesidir; tek koşul sonunda elde edilen fonksiyonun tümel
olmasıdır. Dolayısıyla tarihsel bir kalıntı olan ``genel'' sözcüğü yanıltıcı
bir addır: Yüzeyde \olref{defn:general-recursive},
\olref[par]{defn:recursive-fn} tanımından \emph{daha az} geneldir. Neyse ki
fark yalnızca görünüştedir. Tanımlar farklı olsa da genel özyinelemeli
fonksiyonlar kümesi ile özyinelemeli fonksiyonlar kümesi aynıdır.

